\section{Introdução}
\label{cap:introducao}

Este relatório realiza a análise de diferentes modelos de regressão para o problema de avaliação de imóveis, utilizando o dataset \textit{Real Estate Valuation}. A análise comparativa é estendida para incluir os resultados do trabalho de referência de Yeh \& Hsu (2018), o que estabelece um benchmark para os modelos desenvolvidos. Entre eles estão o de Mínimos Quadrados (MQ), o Perceptron Logístico (PS), o Perceptron Multicamadas (MLP) com uma e duas camadas ocultas (MLP-H1 e MLP-H2), além de modelos de referência provenientes de bibliotecas. 

A análise subsequente mostrará que, quando acompanhado de um pré-processamento criterioso dos dados, até mesmo modelos lineares podem rivalizar com arquiteturas não lineares mais complexas, enquanto as abordagens baseadas em árvores de decisão se destacam por oferecer o desempenho mais robusto diante deste problema. 